\section{The original material generator in the arithmetic coordinate}
\label{sec:material-arithmetic-generator}

We retain the polynomial $F$ and the incompressible vector field $U$ of
Section~\ref{sec:inverse-fibre-heat}, originating in the polynomial
construction cited there\cite{tao2026jacobian,speyer2026jacobian,poplett2026alias}.
The arithmetic functions and their heat parameter are precisely those
of Section~\ref{sec:connes-complete-transport}, with
$\Xi(s)=\xi(1/2+is)$ and $Z_t(s)=8H_t(2s)$
\cite{ccmProlate2024,rodgersTao2020}.
The following calculations give their actual pullback by $F$ and the
original Euclidean diffusion, rather than replacing that diffusion by a
chosen scalar heat operator.

\subsection{An injective arithmetic pullback and the full material drift}

Write $q=(x,y,w)$ and define
\[
 s(q)=F_1(q)/2,\qquad b(q)=F_2(q),\qquad c(q)=F_3(q).
\]
These names denote specific functions on the original space, with no
claim that $F$ has a global inverse. For $W_j=(DF)^{-1}e_j$ and
$U=2W_1+6F_3W_2$, the identity $W_jF_k=\delta_{jk}$ gives
\begin{equation}
 Us=1,\qquad Ub=6c,\qquad Uc=0,\qquad
 U\bigl(\phi(s,b,c)\bigr)=(\partial_s\phi+6c\partial_b\phi)(s,b,c).
 \label{eq:mag-full-drift}
\end{equation}
The chain rule proves the last identity for every differentiable $\phi$,
over $\R$, and every holomorphic $\phi$ over $\C$. In particular
$b-6cs=F_2-3F_3F_1$ is a first integral. This retains the $b$ drift
when $c\ne0$.

For entire functions let
\[
 \mathcal I:\mathcal O(\C)\longrightarrow\mathcal O(\C^3),
 \qquad (\mathcal I f)(q)=f(s(q)).
\]
It is an injective algebra map. Indeed, the polynomial section
$\iota(s)=(0,0,2s)$ satisfies $s(\iota(s))=s$, so
$\iota^*\mathcal I f=f$. Both maps are continuous for uniform convergence
on compact sets, since continuous images of compact sets are compact.
The same section identifies the image with the original function
space when that image is given the transported topology. Repeated
application of \eqref{eq:mag-full-drift} proves
\begin{equation}
 U^k\mathcal I f=\mathcal I f^{(k)}\quad(k\ge0).
 \label{eq:mag-intertwiner}
\end{equation}
Thus this is an actual intertwiner, including its explicit left inverse.

On the order-at-most-one function space of
Section~\ref{sec:cq-heatspace}, the already proved heat group gives
$\mathscr H_t=\mathcal I U_t\iota^*$ on its image. Its series is
\[
 \mathscr H_t\mathcal I f
 =\sum_{k\ge0}\frac{(-t/4)^k}{k!}U^{2k}\mathcal I f
 =\mathcal I(U_tf).
\]
Convergence follows from the proved derivative estimates there, transported
by $\mathcal I$; no convergence assertion on every entire function of three
variables is needed or made. In particular the original arithmetic input,
not a polynomial proxy, obeys
\begin{equation}
 \widetilde Z_t(q)=Z_t(s(q)),\qquad
 \partial_t\widetilde Z_t=-\tfrac14U^2\widetilde Z_t,\qquad
 \widetilde Z_0=\mathcal I\Xi .
 \label{eq:mag-actual-heat}
\end{equation}
If $X_\tau(q)$ is the maximal material flow where it exists, integration
of \eqref{eq:mag-full-drift} gives
\[
 (s,b,c)(X_\tau q)=(s(q)+\tau,b(q)+6c(q)\tau,c(q)).
\]
Hence $(\mathcal I f)(X_\tau q)=f(s(q)+\tau)$ on that domain.
The material parameter $\tau$ translates $s$; the arithmetic heat
parameter $t$ generates $-U^2/4$. Their different actions are related
by these equations, not identified by a change of name.

\subsection{The branch derivative with every coefficient retained}

In the inverse-root space use the full relation
\[
 P(s,b,c;r)=cr^3-2r^2+br-4s=0.
\]
At $P_r\ne0$, differentiation with $b,c$ fixed gives
$\partial_s r=4/P_r$. The lifted material vector field, including its
coefficient drift, is instead
\begin{equation}
 U_{\rm root}=\frac{4-6cr}{P_r}\partial_r+6c\partial_b .
 \label{eq:mag-root-drift}
\end{equation}
Indeed $U_{\rm root}P=(4-6cr)+6cr-4=0$ when the independent $s$
component is included as in \eqref{eq:mag-full-drift}. More explicitly,
on the root hypersurface with local coordinates $(r,b,c)$ one has
$s=(cr^3-2r^2+br)/4$, and the displayed operator gives
$U_{\rm root}s=1$, $U_{\rm root}b=6c$, $U_{\rm root}c=0$.
This proves the stated lift without deleting the transverse drift.

At $c=0$ the coefficient $b$ is constant along the flow. Retain it:
\[
 s=\frac{br}{4}-\frac{r^2}{2},\quad
 \delta_b=b^2-32s=(b-4r)^2,\quad
 \alpha=\frac{b-4r}{2},\quad D_b=\frac4{b-4r}\partial_r.
\]
On a base open set on which $\delta_b\ne0$, the quadratic algebra
$\mathcal O[r]/(r^2-br/2+2s)$ is free in the ordered basis $(1,r)$.
Its multiplication and connection matrices are
\begin{equation}
 R_b=\begin{pmatrix}0&-2s\\1&b/2\end{pmatrix},\qquad
 D_b=\partial_s I+A_b,\qquad
 A_b=\begin{pmatrix}0&4b/\delta_b\\0&-16/\delta_b\end{pmatrix}.
 \label{eq:mag-connection-b}
\end{equation}
To check the second formula, use
$D_br=4/(b-4r)=(4b-16r)/\delta_b$ and differentiate
$f_0(s)+rf_1(s)$ by the product rule. In particular
\begin{align}
 R_b^2-\tfrac b2R_b+2sI&=0,\notag\\
 [D_b,R_b]&=4(bI-4R_b)/\delta_b,\notag\\
 D_b^2&=\partial_s^2 I+2A_b\partial_s+(16/\delta_b)A_b.
 \label{eq:mag-connection-b-square}
\end{align}
For the commutator calculate $R_b'+A_bR_b-R_bA_b$.
For the square use $A_b'=(32/\delta_b)A_b$ and
$A_b^2=-(16/\delta_b)A_b$. These computations prove every
entry, including the zeroth-order term. At $b=0$ the connection is
$\operatorname{diag}(\partial_s,\partial_s+1/(2s))$, the earlier
two-sheet connection with exactly the same constants.

The sheet involution $r\mapsto b/2-r$ has matrix
$J_b=\left(\begin{smallmatrix}1&b/2\\0&-1\end{smallmatrix}\right)$,
with $J_b^2=I$ and $J_bR_bJ_b=bI/2-R_b$.
The algebra trace is
\[
 \operatorname{tr}(f_0+rf_1)=2f_0+\tfrac b2 f_1,\qquad
 \ker\operatorname{tr}=\{(-bh/4,h):h\in\mathcal O\}.
\]
This proves its kernel as well as the retained sheet symmetry.

The excluded locus in this inverse chart is
$r=b/4$, $s=s_*:=b^2/32$, $\alpha=0$.
Writing \emph{explicitly} $\eta=r-b/4$ gives
\begin{equation}
 s=s_*-\eta^2/2,\qquad
 H=\left(-\frac1{2\eta},\frac b4+3\eta,
                  26\eta^2+\frac{3b\eta}{2}\right).
 \label{eq:mag-shifted-branch}
\end{equation}
This is a variable substitution with the inverse $r=\eta+b/4$,
not deletion of $b$. The separate point with $x=0$ remains
$(0,b,2s-4b^2)$, including at $s=s_*$. For complex $b$ every
complex value of $s_*$ occurs; for real $b$ the values are exactly
$[0,\infty)$. Thus zero-freeness at the single $b=0$ branch origin
does not exclude arithmetic zeros at the other branch values.

For precision, if the actual function at a fixed time has
$Z_t(s)=(s-s_*)^m A(s)$ near $s_*$, where $A(s_*)\ne0$, then
substitution proves the complete local identity
\[
 Z_t(s_*-\eta^2/2)=(-1/2)^m\eta^{2m}A(s_*-\eta^2/2).
\]
The ramification doubles that existing multiplicity; it does not
supply a new zero. For $Z_t(s_*)\ne0$ the same pullback is nonzero
at $\eta=0$. Both cases follow directly by evaluating the analytic
factor. These are local formulas for the actual function, not
an assumption about the location or multiplicity of its zeros.

\subsection{The original Euclidean diffusion and its fibre dependence}

For real $q$, put $v=1+xy$. Direct differentiation of the retained
polynomial gives
\begin{align*}
 s_x&=\tfrac12\{3yv^2w+y^3(1+6v)\},\\
 s_y&=\tfrac12\{3xv^2w+2y(v+3v^2)+xy^2(1+6v)\},\\
 s_w&=\tfrac12v^3,\\
 \Delta_Es&=3wv(x^2+y^2)+18x^2y^2+21xy+3y^4+4.
\end{align*}
For completeness, use $F_1=v^3w+y^2(v+3v^2)$, with
$v_x=y$, $v_y=x$, $v_w=0$, to obtain the three first derivatives.
The second $x$ derivative is $3y^2vw+3y^4$.
The second $y$ derivative is
$3x^2vw+v+3v^2+2xy(1+6v)+3x^2y^2$.
Their sum, with $s_{ww}=0$ and $v=1+xy$, gives the displayed
Laplacian. The scalar chain rule is therefore exactly
\begin{equation}
 \Delta_E(\mathcal I f)
 =|\nabla_Es|^2\,\mathcal I f''+(\Delta_Es)\,\mathcal I f'.
 \label{eq:mag-Euclidean-diffusion}
\end{equation}
The squared gradient here is the Euclidean real squared norm. For
holomorphic algebraic extension its expression is $\sum_j s_{q_j}^2$,
not a silently inserted Hermitian norm.

Along the actual escaping trajectory
\[
 q(z)=(z^{-1},-3z/2,13z^2/2),\qquad
 z=\sqrt{1-8\tau}>0,
\]
substitution gives $v=-1/2$, $(s,b,c)=(-z^2/8,0,0)$ and
\begin{equation}
 \nabla_Es=\left(-\frac{9z^3}{32},-\frac{3z}{16},-\frac1{16}\right),
 \quad |\nabla_Es|^2=\frac{81z^6+36z^2+4}{1024},
 \quad \Delta_Es=\frac{13-27z^4}{4}.
 \label{eq:mag-orbit-diffusion}
\end{equation}
For example the $x$ component is
$\{-117z^3/16+27z^3/4\}/2=-9z^3/32$,
by direct substitution into the preceding derivative formula.
The squared norm is the sum of the three
displayed squares; the Laplacian follows from its polynomial formula.
Consequently
\[
 (\Delta_E\mathcal I f)(q(z))
 =\frac{81z^6+36z^2+4}{1024}f''(-z^2/8)
       +\frac{13-27z^4}{4}f'(-z^2/8).
\]
Both coefficients have finite limits, $1/256$ and $13/4$, as
$z\downarrow0$. Meanwhile $q(z)$ leaves every bounded set.
This is a statement about the specified arithmetic observable along
that trajectory, not a boundedness statement for every transported
covector or for the full material metric.

There is an exact, finite-point reason that the full diffusion cannot
be encoded by a scalar operator depending only on $s$. The points
\[
 q_1=(1,-3/2,13/2),\qquad q_0=(0,0,-1/4)
\]
both have $F(q_j)=(-1/4,0,0)$ and hence $s=-1/8$.
At $q_1$, $\Delta_Es=-7/2$; at $q_0$, $\Delta_Es=4$.
Thus $\Delta_Es$ is not constant even on this full $F$ fibre.
If an operator $L$ on scalar functions of $s$ satisfied
$\Delta_E\mathcal I=\mathcal I L$, applying it to $f(s)=s$
would give the same value at these two points, contradicting
$-7/2\ne4$. The actual relation is instead the completely specified
two-coefficient formula \eqref{eq:mag-Euclidean-diffusion}.
On the escaping trajectory one can write its coefficients as
$-81s^3/2-9s/32+1/256$ and $13/4-432s^2$, respectively,
but that trajectory formula is not an autonomous reduction valid
on every fibre.

\subsection{Transport of the complete arithmetic form}

To carry the same heat map into arithmetic tests, retain
$\mathcal E=C_c^\infty(\R;\C)$ and set, for real $t$,
\[
 G_tf(u)=e^{tu^2/4}f(u),\qquad
 W_t(f,g)=W(G_{-t}f,G_{-t}g).
\]
Multiplication by the smooth nonvanishing function and its inverse
preserves each compact support and its smooth topology. Thus $G_t$
is an automorphism of the test space. Section~\ref{sec:cq-tests}
proved $\mathcal F G_t=U_t\mathcal F$ and
$T_t=G_tTG_{-t}=-i\partial_u+itu/2$ for the retained Fourier sign.
The symmetry of $T$ for $W$ now gives, directly,
\[
 W_t(T_tf,g)
 =W(TG_{-t}f,G_{-t}g)
 =W(G_{-t}f,TG_{-t}g)=W_t(f,T_tg).
\]
Moreover, for any finite family of tests,
\begin{equation}
 [W_t(G_tf_i,G_tf_j)]_{ij}=[W(f_i,f_j)]_{ij}.
 \label{eq:mag-Weil-congruence}
\end{equation}
This is equality of every matrix entry, including its entire quadratic
form and inertia. In particular transport of both tests and pairing
does not manufacture a negative value.

The exact arithmetic expression for fixed, untransported tests uses
\[
 h_t(x)=\int_\R g(x+v)\overline{f(v)}
              e^{-t((x+v)^2+v^2)/4}\,dv
       =(G_{-t}g)*(G_{-t}f)^\star(x).
\]
The convolution identity follows by substituting the definition of
reflection-conjugation. It keeps support inside the fixed difference
of the two supports. Applying the full explicit formula
\cite{weil1952,alpogeFurman2026} yields
\begin{align}
 W_t(f,g)={}&\widehat h_t(i/2)+\widehat h_t(-i/2)\notag\\
 &+\int_\R\widehat h_t(\lambda)
   \left\{\frac1{2\pi}\Re\psi\left(\frac14+\frac{i\lambda}{2}\right)
                       -\frac{\log\pi}{2\pi}\right\}d\lambda\notag\\
 &-\sum_{n\ge2}\frac{\Lambda(n)}{\sqrt n}
             \{h_t(\log n)+h_t(-\log n)\}.
 \label{eq:mag-full-transported-Weil}
\end{align}
The prime-power sum is finite by compact support. Uniform bounds on
all derivatives on a fixed compact $t$ interval give rapid Fourier
decay in horizontal strips by integration by parts; this justifies
both the integral and its parameter derivatives. Differentiating
$G_{-t}$, or equivalently the displayed compact integral, therefore gives
\[
 \frac d{dt}W_t(f,g)
 =-\tfrac14\{W_t(u^2f,g)+W_t(f,u^2g)\}.
\]
Every pole, gamma and prime term is still present in
\eqref{eq:mag-full-transported-Weil}. This is the congruent transport
of the original Weil form; a form defined instead by the moving
zero divisor of $Z_t$ has not been substituted for it. The exact
transported quotient and the moving divisor comparison were proved
in Section~\ref{sec:connes-complete-transport}. The present
calculation supplies the actual material generator, the full
Euclidean diffusion and the retained arithmetic test map. It
does not exhibit a negative value of the original Weil form.
